% !TeX root = main.tex
% =============================================================================
% SEÇÃO 2: ATIVIDADES REALIZADAS
% =============================================================================
% Esta seção deve apresentar:
% - Um relato detalhado de todas as atividades executadas durante o período
% - A metodologia utilizada em cada atividade
% - Os resultados obtidos até o momento
% - Ferramentas, técnicas e recursos utilizados
% - Contribuições científicas (publicações, submissões, apresentações)
% =============================================================================

\section{Atividades Realizadas}

% Introdução geral sobre o período e as atividades desenvolvidas
A seguir, estão sintetizadas as principais atividades realizadas durante o 
período de [mês/ano] a [mês/ano]:

\begin{enumerate}
    % --- Atividade 1 ---
    \item [\textbf{i.}] \textbf{[Título da Primeira Atividade]:} 
    [Descreva detalhadamente a primeira atividade realizada, incluindo 
    objetivos, metodologia aplicada, ferramentas utilizadas e resultados 
    obtidos. Cite referências relevantes quando apropriado. Para citar, use \cite{autor2022_dissert} ou \citeonline{autor2022_dissert} conforme o estilo de citação adotado.] 

    
    % --- Atividade 2 ---
    \item [\textbf{ii.}] \textbf{[Título da Segunda Atividade]:} 
    [Descreva a segunda atividade, destacando os principais aspectos 
    metodológicos e os resultados alcançados. Se utilizou ferramentas 
    específicas, mencione-as e justifique sua escolha.\footnote{Você pode 
    adicionar notas de rodapé para informações complementares, como links 
    para ferramentas ou recursos utilizados.}]
    
    % Exemplo: A ferramenta selecionada foi baseada em critérios estabelecidos
    % por \citeonline{organizacao2024_manual}, demonstrando eficácia comprovada
    % em estudos anteriores \cite{autor2024_conf,autor2023_exemplo}.
    
    % --- Atividade 3 ---
    \item [\textbf{iii.}] \textbf{[Título da Terceira Atividade]:} 
    [Continue descrevendo as atividades realizadas, mantendo o padrão de 
    detalhamento e clareza. Inclua visualizações, gráficos ou tabelas 
    quando relevante para ilustrar os resultados.]
    
    % --- Atividade 4 (Opcional) ---
    % Adicione mais atividades conforme necessário
    \item [\textbf{iv.}] \textbf{[Submissões e Contribuições Científicas]:} 
    [Se houver submissões de artigos, participação em eventos, ou outras 
    contribuições científicas, descreva-as neste item. Inclua informações 
    como nome do evento/periódico, status da submissão, e principais 
    contribuições apresentadas.]
    
    % --- Atividade 5 (Opcional) ---
    \item [\textbf{v.}] \textbf{[Aspectos Éticos e Regulatórios]:} 
    [Se seu projeto envolver pesquisa com seres humanos ou outros aspectos 
    que requeiram aprovação ética, descreva as ações realizadas nesse 
    sentido, como submissão e aprovação pelo Comitê de Ética em Pesquisa.]
    
\end{enumerate}

% Dica: Você pode adicionar subseções se alguma atividade for muito extensa
% ou necessitar de maior detalhamento.